\section{Research}
\label{sec:background}
% Rubric: exposition that synthesises the literature, in your own words.
% Budget: ~800 words.




% General Explanation of Reaction Diffusion System
This project explores a system of partial differential equations known as Reaction Diffusion Systems. Reaction Diffusion Systems emerge from chemistry and are used as a tool to model how substances from a chemical reaction *diffuse* in space. Hence, the name Reaction Diffusion Systems.

% General Form of a Reaction Diffusion System: 
\begin{align*}
	\partial_t q = D \nabla^2 q + R(q)
\end{align*}


% Research from complexity explorer (youtube video) 

% The following is an example system shown in the video:
Reaction of two chemicals $U$ and $V$ that follow these rules:
\begin{align*}
	U + 2V & \to 3V \\
	V      & \to P
\end{align*}

$V \to P$ represents the rate of decay of chemical $V$


Thus, the reaction diffusion system based on the above rule is given by:
\begin{align*}
	\frac{\partial U}{\partial t}  = - UV^2 + F(1 - U) + D_u (\Delta) \\
	\frac{\partial V}{\partial t}  = UV^2 + (F - K)V  + D_v (\Delta)
\end{align*}

Graphing the reaction diffusion system produces an interesting pattern? But i don't understand how this example works?









